\chapter*{Trang thông tin luận văn}
\addcontentsline{toc}{chapter}{Trang thông tin luận văn tiếng Việt}
\begin{center}\textbf{\large TRANG THÔNG TIN LUẬN VĂN}\end{center}

\renewcommand{\arraystretch}{1.25}
\noindent\begin{tabularx}{\textwidth}{@{}lX@{}}
Tên đề tài luận văn: & \TenDeTaiMotDong. \\
Ngành: & \Nganh \\
Mã số ngành: & \MaNganh \\
Họ và tên HVCH: & \HocVien \\
Khóa đào tạo: & \Khoa \\
GVHD: & \GVHD \\
Cơ sở đào tạo: & Trường ĐH Khoa học Tự nhiên, ĐHQG-HCM \\
\end{tabularx}

\subsection*{1. Tóm tắt nội dung luận văn}

Các mô hình giao diện hiện nay sinh ra toạ độ để máy tự bấm. Luận văn nghiên cứu một
đầu ra khác trên cùng dữ liệu vào: một câu tiếng Anh nói cho \emph{người đọc} biết phải
chạm vào phần tử nào trên màn hình. Trở ngại chính không nằm ở khâu sinh câu mà ở khâu
chấm: cùng một nút có nhiều cách gọi tên hợp lệ, nên không câu chuẩn đơn lẻ nào gánh
được vai trò đáp án.

Luận văn đề xuất thước đo \textbf{executability}: câu sinh ra được đưa cho một mô hình
định vị phần tử độc lập, mô hình này không hề thấy toạ độ đáp án; bước được tính đúng
khi điểm nó trả về rơi vào ô Voronoi của phần tử người dùng đã chạm và loại thao tác
được gọi tên đúng. Thước được kiểm bằng bốn phép độc lập trước khi dùng để kết luận:
một cổng chặn về sai số của dụng cụ đo (đạt: sai số trung vị $0{,}7\%$ bề ngang màn so
với ngưỡng $3\%$ khoá trước), một phép chọn luật trúng dựa trên \emph{sàn} đo được (luật
dung sai quy ước cho điểm $84{,}3\%$ số câu chỉ sai nút, luật ô Voronoi chỉ $2{,}8\%$),
một phép đo trần ($75{,}7\%$, KTC 95\% $[74{,}1; 77{,}3]$ khi đưa chính câu người viết đi
chấm), và một bộ bơm lỗi với ngưỡng đóng băng trước khi chạy.

Về phía mô hình, luận văn xây dựng hoàn toàn tự động một tập giám sát $64.567$ bước từ
bộ AndroidControl, không thuê người dán nhãn, đã kiểm rò rỉ dạy--kiểm ở quy mô đủ và
không tìm thấy tác vụ nào trùng. Thành phần đề xuất là ``mô tả phân biệt trước, phát
ngôn sau'': đích huấn luyện chứa một dòng khai báo có cấu trúc về phần tử đích trước khi
tới câu, và dòng khai báo bị cắt bỏ trước khi chấm.

Trên $4.463$ bước chạm của tập kiểm, mô hình Qwen2.5-VL-3B tinh chỉnh QLoRA đạt
$59{,}1\%$ so với $47{,}6\%$ của chính mô hình đó khi chưa tinh chỉnh; chênh lệch ghép
cặp $11{,}5$ điểm ($p<0{,}001$) đứng vững qua năm cách giải thích thay thế, còn cách thứ
sáu (bắt chước văn phong) chưa loại trừ được.

\subsection*{2. Những kết quả mới của luận văn}

\begin{itemize}[leftmargin=1.4em]
\item \textbf{Thước đo không cần đáp án tham chiếu cho bài toán sinh hướng dẫn giao
diện.} Luật trúng được chọn dựa trên hành vi đo được với đầu vào sai nút, không dựa
trên ngưỡng dung sai quy ước; thứ tự giữa các hệ thống giữ nguyên dưới năm luật chấm
khác nhau.

\item \textbf{Trần của thước được đo, không phải giả định.} Đưa chính câu do người viết
đi chấm cho $75{,}7\%$; mọi điểm số của mô hình phải đọc trên nền này chứ không phải
trên nền $100$. Ước lượng cũ $70{,}0\%$ dựa trên mẫu con $300$ bước lệch $5{,}7$ điểm và
đã bị rút.

\item \textbf{Đường ống dữ liệu tự động.} Chuyển AndroidControl thành tập giám sát cho
tác vụ mới: ghép hai bản phân phối trên Hugging Face, kiểm phép ghép bằng đối chứng lệch
một bước ($48\%$ so với $20\%$), dựng nhãn mô tả từ cây trợ năng kết hợp OCR, chín bất
biến cấu trúc đều đạt ở quy mô đủ.

\item \textbf{Bản đăng ký trước có luật đọc cho cả bốn kết cục}, đưa vào kho phiên bản
trước lượt huấn luyện đầu tiên; mọi thay đổi về sau ghi thành mục sửa đổi có ngày.

\item \textbf{Định vị bài toán.} Trong bộ AndroidControl gốc và trong các công trình
dùng bộ này, câu hướng dẫn từng bước do người viết nằm ở \emph{đầu vào}; ở đây nó là
\emph{đích sinh} và là đầu ra duy nhất đem chấm.
\end{itemize}

\vspace{1.5em}
\noindent\begin{tabularx}{\textwidth}{@{}XX@{}}
\centering CÁN BỘ HƯỚNG DẪN & \centering\arraybackslash HỌC VIÊN CAO HỌC \\
\centering (Ký tên, họ tên) & \centering\arraybackslash (Ký tên, họ tên) \\
\end{tabularx}

\vspace{4em}
\begin{center}
XÁC NHẬN CỦA CƠ SỞ ĐÀO TẠO\\ HIỆU TRƯỞNG
\end{center}
\clearpage

% ==================== BẢN TIẾNG ANH ====================
\chapter*{Thesis information}
\addcontentsline{toc}{chapter}{Trang thông tin luận văn tiếng Anh}
\begin{center}\textbf{\large THESIS INFORMATION}\end{center}

\noindent\begin{tabularx}{\textwidth}{@{}lX@{}}
Thesis title: & \TenDeTaiEN. \\
Speciality: & Computer Science \\
Code: & \MaNganh \\
Name: & Le Doan Phuong Uyen \\
Academic year: & \Khoa \\
Supervisor: & \GVHDen \\
At: & VNUHCM -- University of Science \\
\end{tabularx}

\subsection*{1. Summary}

GUI agents produce coordinates for a machine to click. This thesis studies a different
output over the same input: a short instruction telling a \emph{person} which on-screen
element to touch. The obstacle is not generation but evaluation, since many wordings
identify the same button and no single reference can carry the judgement.

We propose \textbf{executability}: the generated sentence is passed to a separate GUI
grounding model that never sees the gold coordinate, and a step counts as correct when
the returned point falls in the Voronoi cell of the element the user touched and the
described action type matches. The instrument is characterised before it is trusted: an
adequacy gate on grounder error (passed, median 0.7\% of screen width against a
threshold of 3\% fixed in advance), a hit rule selected from its measured floor (the
conventional tolerance test credits a sentence naming the wrong button in 84.3\% of
cases, the cell rule in 2.8\%), a measured ceiling of 75.7\% (95\% CI
$[74.1, 77.3]$) obtained by scoring the human-written references themselves, and a
perturbation self-check with thresholds fixed before execution.

On the model side we build 64,567 supervision steps from AndroidControl fully
automatically, with zero task-level train--test leakage verified at full scale. The
proposed component is descriptor-first generation: the training target carries a
structured descriptor of the target element before the sentence, and the descriptor is
removed before scoring.

Over 4,463 held-out touch steps a QLoRA-tuned Qwen2.5-VL-3B reaches 59.1\% against
47.6\% for the same model untuned; the paired gap of 11.5 points ($p<0.001$)
survives five checks for alternative explanations; a sixth, imitation of the annotator
register, is not ruled out.

\subsection*{2. Novelty of thesis}

\begin{itemize}[leftmargin=1.4em]
\item A reference-free metric for element identification in generated GUI instructions,
with the hit rule selected from measured behaviour on wrong-element inputs rather than
from a tolerance convention, and shown to preserve system ordering under five different
rules.
\item A measured, not assumed, ceiling: human-written references score 75.7\%, the
level every model score must be read against. An earlier 70.0\% estimate from a
300-step subsample was five points low and has been retracted.
\item A fully automatic pipeline converting AndroidControl into supervision for this
task, with the cross-modal join validated against a deliberately shifted control
(48\% against 20\%).
\item A pre-registered ablation design fixing the reading of all four outcomes,
committed to version control before the first training run.
\item Positioning: in AndroidControl and in work built on it, the human-written
per-step instruction is an \emph{input}; here it is the generation target and the only
scored output.
\end{itemize}

\vspace{1.5em}
\noindent\begin{tabularx}{\textwidth}{@{}XX@{}}
\centering SUPERVISOR & \centering\arraybackslash MASTER STUDENT \\
\end{tabularx}

\vspace{4em}
\begin{center}
CERTIFICATION\\ UNIVERSITY OF SCIENCE\\ PRESIDENT
\end{center}
\clearpage
